\chapter{Theoretical Foundations and Literature Review}
\label{ch:foundations}

This chapter sets out the bodies of theory the thesis relies on --- latency
arbitrage and discrete-time market design, adverse selection, market
transparency, uniform-price auction theory and the measurement of market
quality. The thesis's own argument that a predictable batch close turns
price improvement into a boundary repricing race builds on these strands,
but is a consequence of a design choice and is therefore developed in the
design chapter (\Cref{sec:priceimprovement}). Each strand is used later
in one of two ways: it shapes a design block of \Cref{ch:design}, or it
motivates a metric of the empirical comparison and predicts its sign. Each
section states that link explicitly, and \Cref{sec:theorylinks} collects
them in one table. Notation and the formal definition of every metric are
in \Cref{ch:background}.

\section{Discrete-Time Market Design}
\label{sec:discretetime}

\subsection{The Mechanics of Latency Arbitrage}

The argument of \citet{budish2015hft} rests on a single structural
observation, not on irrationality or market power: in a continuous
market, a public information event creates a discrete arbitrage
opportunity that is allocated by relative speed.

Consider a liquidity provider quoting a two-sided market in an asset
whose fair value is tracked by a correlated, continuously observable
signal. When the signal jumps, every participant learns the same thing at
essentially the same instant: the liquidity provider wants to cancel its
stale quote, every other fast participant wants to hit it. The engine
processes both messages in arrival order, so whether the quote is
cancelled or picked off depends on network path length, hardware and code
efficiency, not on anything informational.

Two consequences follow. The opportunity is mechanically recurrent,
appearing every time the public signal moves --- thousands of times a day
in a liquid instrument. And the value of being marginally faster does not
decline as everybody invests in speed: since only the winner is paid, the
equilibrium is a rent-dissipating race in which expenditure on latency is
bounded by the value of the opportunities, which persist however fast the
market becomes. Speed is a purely relative good.

The cost is not borne by the racers. A liquidity provider that expects to
lose a fraction of its volume to picked-off quotes must recover that loss
from the trades it does not lose, shifting the burden onto liquidity
demanders as a wider spread --- a tax on liquidity in the words of
\citet{budish2015hft}.

\subsection{Empirical Magnitude}

\citet{aquilina2022quantifying} measure this directly, using London Stock
Exchange message-level data (including messages never executed) to
identify latency races as clusters of near-simultaneous attempts to take
or cancel the same quote. Races are extremely frequent and short --- far
shorter than any batch interval measured in milliseconds; a large share of
the volume at the exact moment of a price change is race volume; and the
aggregate transfer, small relative to total volume but large in absolute
terms, concentrates in the effective spread paid by liquidity demanders.
Removing the race therefore removes a component of the spread. This is why
the effective spread, and its split into realized spread and price impact
(\Cref{sec:metrics21}), is the central cost measure of RQ2.1.

\subsection{Batching as the Response}

Frequent batch auctions change two rules at once, and the two do
different work. \emph{Discrete time} removes the informational value of
within-interval arrival order: if orders are collected over $[t,
t+\tau)$ and matched only at $t+\tau$, any two orders inside the window
are symmetric with respect to time as far as the price is concerned
(whether they stay symmetric in allocation depends on the rationing rule,
\Cref{sec:rationing}). \emph{Uniform pricing} removes the
value of queue position within the matched set: if all executions occur
at one price $\clr$, no participant obtains a better price by being
earlier. Either rule alone is insufficient --- discrete time with
discriminatory pricing would still reward earlier submission, and uniform
pricing in continuous time is not well defined. Together they imply that
speed pays only for crossing an interval boundary, getting into the
current batch rather than the next: the value of a marginal nanosecond
falls from ``wins the race'' to ``matters only in the vanishingly rare
case that it changes which batch an order lands in''. In the model of
\citet{budish2015hft}, liquidity providers then no longer need to price
sniping risk into their quotes and spreads narrow. The residual cost is
delay, an average of $\tau/2$ if orders arrive uniformly, so the choice of
$\tau$ trades speed neutralisation against delay imposed on everybody.
The benefit is not free of side effects: \citet{baldauf2020hft} show that
competition among fast traders affects both liquidity and the incentive to
acquire information, so a measure that curbs the race can lower spreads
and still change how much information prices contain.
This conclusion assumes that waiting inside the window gains nothing;
\Cref{sec:priceimprovement} shows that with a predictable close and a
visible batch it does, and that speed then regains value near the
boundary.

Both halves of this trade-off are followed up later. The choice of $\tau$
and of its closing rule is design Block~1 (\Cref{sec:interval}). The
benefit, a lower adverse-selection cost, is measured by the spread
decomposition (RQ2.1), and the cost, delay, by the time to execution
(RQ2.3, \Cref{sec:metrics23}).

\section{Information Asymmetry and Adverse Selection}
\label{sec:adverse}

Latency arbitrage is one specific mechanism through which a liquidity
provider loses money to better-informed counterparties; the general
theory of that loss predates high-frequency trading by three decades and
is what makes the spread a meaningful economic object.

\subsection{The Spread as an Information Premium}

\citet{glosten1985bid} model a competitive, risk-neutral market maker
facing a mixed population of informed and uninformed traders it cannot
tell apart. Because a buy order is more likely to come from an informed
trader when the asset is undervalued, every order carries information, so
quotes must be conditionally fair given the direction of the order
expected --- a bid--ask spread that exists purely as an information
premium, even with no inventory risk or processing cost. The spread widens
with the proportion of informed traders and the precision of their
information, and narrows when the market maker can identify a subset of
flow as uninformed. This is the theoretical foundation of order flow
segregation: separating benign flow lets a liquidity provider quote
tighter against it without cross-subsidising losses on the rest. It
underlies design Block~3 (\Cref{sec:segregation}). The simulated engine
clears a single pool (\Cref{tab:fbachoices}), so segregation is analysed
in the design part but not measured.

\subsection{Strategic Informed Trading and Market Depth}

\citet{kyle1985continuous} approaches the same problem from the informed
trader's side: a single informed agent chooses how much to trade, knowing
that trading more reveals more, while noise traders provide cover and a
competitive market maker sets the price as the conditional expectation of
value given total order flow. The equilibrium price function is linear,
$p = \mu + \lambda f$, where $\mu$ is the prior expected value, $f$ is
aggregate order flow and
$\lambda$ measures price impact per unit of flow; its reciprocal, market
depth, is the volume required to move the price by one unit. Kyle's model
is a batch model --- the market maker observes aggregate flow and sets a
single price, the structure of a uniform-price call auction --- so the FBA
literature can be read as asking what happens when a real market is
organised the way Kyle already assumes. $\lambda$ also supplies the
empirical bridge: price impact is the operational measure of adverse
selection. The thesis estimates $\lambda$ for each engine as the slope of
price change on order flow (\Cref{sec:formularq22}); a lower $\lambda$
would mean prices react less to one-sided flow. Both engines are regressed
on the same measure of flow, the venue's own net taker flow, so that the
two slopes can be compared directly.

\subsection{Decomposing the Spread}
\label{sec:spread-decomposition}

The adverse selection premium is not directly observable; only the total
cost of a trade is. The standard strategy is to decompose the effective
spread into a realized spread and a price impact component
\citep{hasbrouck1993assessing,hendershott2011algorithmic}, as defined in
\eqref{eq:effspread}--\eqref{eq:priceimpact}, with $S^{\mathrm{eff}}_k =
S^{\mathrm{real}}_k(\Delta) + \mathrm{PI}_k(\Delta)$. The effective spread
is what the liquidity demander paid relative to the midpoint; the realized spread is what the provider still holds once the
market has revised its valuation, the gross compensation for supplying
liquidity; the price impact is the permanent revision of the midpoint in the
direction of the trade, the part lost to the information the order carried.

When a public signal moves, discrete uniform-price clearing removes the
latency-arbitrage mechanism, since the aggressive order and the maker's
cancellation are processed together and execution occurs at a price that
already reflects orders submitted in response to the signal. It does not
remove adverse selection in the theoretical sense --- a genuinely
informed trader can still price aggressively into a batch
\citep{eibelshauser2022fba} --- so the prediction is a reduction in
price impact, not its elimination. A fall in effective spread driven by
price impact is therefore the result theory predicts; a fall driven by
realized spread would mean providers simply earn less with the
informational component untouched, and a fall in price impact with no
change in effective spread would mean the benefit is captured by
providers rather than passed on. These three readings are what the
RQ2.1 comparison in \Cref{sec:res21} is designed to tell apart, which is
why all three components are reported rather than the effective spread
alone.

\section{Market Transparency Theory}
\label{sec:transparency}

\citet{madhavan1996transparency} defines transparency as participants'
ability to observe information about the trading process, and shows more
of it is not unambiguously better: greater pre-trade transparency helps
prices aggregate information and reduces uncertainty about market state,
but also exposes the trading intent of participants transacting in size,
and can increase volatility and reduce liquidity in some configurations.
Consolidating with \citet{ohara1995microstructure}, disclosure design is
a genuine trade-off, not a monotone improvement.

In an FBA the question is narrower: what is disclosed during the
accumulation phase, before the auction clears? Moving toward opacity
protects submitted orders from being exploited but degrades the
information available for pricing, and classical auction theory predicts
that participants respond to uncertainty about the clearing price by
shading their bids \citep{klemperer1999auction}. In a uniform-price
setting shading is not merely a transfer: shaded limit prices are less
likely to cross, so it can reduce the volume that clears.

The trade-off assumes the operator can choose where to sit among these
options. Decentralised blockchain venues are the limiting case where it
cannot: a transaction waiting in a public mempool is visible to everyone
before it executes, so such markets are fully open by default without the
benefit that regime is supposed to buy. Disclosure normally exposes the trading intent of the order submitter
\citep{madhavan1996transparency}, and the informational disadvantage this
creates normally falls on the market maker posting a quote
\citep{glosten1985bid}. A mempool flips both of these. What is exposed is
not a resting quote but an instruction that has not executed yet, and the
party exposed by it is not the quoter but the one demanding liquidity,
the submitter itself. An observer can buy ahead of a pending swap and
sell into the price it moves, sandwiching the submitter
\citep{daian2020flash}. Batch
clearing removes this by construction: if every order in the interval
settles at one uniform price, position within the batch carries no value
for the price an order obtains.

This strand shapes design Block~2 (\Cref{sec:disclosure}). It cannot be
tested in the replay: the replayed orders never see what the engine
publishes, so the disclosure regime cannot change any result
(\Cref{tab:fbachoices}), and bid shading is listed among the behavioural
predictions out of reach (\Cref{sec:limitations}).

\section{Auction Theory and Uniform-Price Clearing}
\label{sec:auctiontheory}

In a uniform-price auction all winning participants transact at the same
market-clearing price, independent of the prices they submitted. A
participant's own bid determines \emph{whether} it trades but, except for
the marginal order, not \emph{at what price}. Misreporting has no upside:
bidding below one's valuation only risks losing an acceptable trade,
bidding above only risks winning an unacceptable one, so for a
price-taking participant with single-unit demand truthful revelation is
approximately optimal, a property inherited from
\citet{vickrey1961counterspeculation}. This is the formal basis for the
claim that batching substitutes competition in price for competition in
speed: if submitted limits are approximately truthful, the clearing price
aggregates valuations rather than reaction times. For a price-taking
participant, visibility of the accumulated batch does not disturb this ---
a participant who can anticipate the clearing price pays it regardless of
the limit it submits above it --- but it does change \emph{when} orders
are submitted, since
observing the indicative price gives a participant reason to withhold
until shortly before the batch closes.

The qualification carried by ``approximately'' depends jointly on
pivotality and size. A participant too small to move the crossing point
controls only its own inclusion, so shading saves nothing and truthful
bidding is optimal without qualification. A pivotal participant with
single-unit demand does pin the clearing price with its own limit, but the
marginal unit carries almost no surplus, so the incentive to shade is
negligible. The case that matters is the pivotal participant with
multi-unit demand: its last unit sets the price for every unit it buys, so
shading the marginal bid reduces the cost of the entire inframarginal
block while risking only the one unit barely worth trading. The saving
scales with size, the loss does not, and the incentive to misreport is
restored; \citet{klemperer1999auction} surveys this literature. Pivotality
is relative to the liquidity accumulated within one batch, not the market
as a whole, so the strength of the truthfulness property is itself a
function of interval length: shorter intervals produce thinner batches,
where an order that would be atomistic in a long window can be
price-setting in a short one.

For every executed non-marginal order, uniform pricing also produces price
improvement relative to the submitted limit: a buyer willing to pay
$\pi_k$ who clears at $\clr < \pi_k$ captures $(\pi_k - \clr) x_k$ of
surplus that, under continuous discriminatory execution at the resting
price, would in general have accrued differently. Aggregating this across
executed orders gives the realized trader surplus metric
(\Cref{sec:metrics23}). \Cref{sec:priceimprovement} shows that for a given
set of executed orders the total of this surplus does not depend on the
price: the price only decides how it is split between buyers and sellers,
so price improvement on one side is a transfer from the other, not a gain.

These results inform design Block~4, the clearing and rationing rule
(\Cref{sec:clearing}). The link between interval length and pivotality
also has an empirical counterpart: thin batches show up as batches that
do not cross and as a high unexecuted residual share (\Cref{sec:metrics23}).

\section{Measurement of Market Quality: Motivating the Metrics}
\label{sec:quality}

The empirical comparison relies on the standard microstructure toolkit,
and each metric of \Cref{sec:metricdefs} is chosen because one of the
theories above makes a prediction about it.

\textbf{Liquidity and transaction costs (RQ2.1).} Quoted spread and depth
describe the cost a hypothetical trade would face: if batching pools
liquidity, as its proponents argue, the FBA should be at least as deep.
Effective spread describes what an actual trade cost, and realized spread
and price impact decompose that cost into the part retained by the
liquidity provider and the part lost to adverse selection
\citep{hasbrouck1993assessing}; \Cref{sec:spread-decomposition} gives the
prediction for each. \citet{amihud2002illiquidity} provides a
low-frequency summary of price impact per unit of volume, robust and easy
to compare across mechanisms because it needs neither a book nor a trade
direction.

\textbf{Price discovery (RQ2.2).} Batching aggregates orders before
setting a price, so it should change how much prices move and how they
respond to order flow. Kyle's lambda (\Cref{sec:adverse}) measures the
price response to flow directly, intra-interval dispersion checks the
mechanical promise of uniform pricing, and volatility measures how much
the price varies. \citet{hasbrouck1995onesecurity} supplies a framework
for attributing price discovery across venues, but it needs venues that
trade the same asset simultaneously and influence each other; the two
engines here never interact, so the thesis compares their price series on
a common time grid instead (matched-timestamp volatility).

\textbf{Execution and timing (RQ2.3).} The cost of batching is delay
(\Cref{sec:discretetime}), measured by the time to execution. Uniform
pricing promises price improvement (\Cref{sec:auctiontheory}), measured by
trader surplus. The boundary argument of the design chapter
(\Cref{sec:priceimprovement}) predicts that, under a fixed and visible
close, orders bunch near the batch close, measured by boundary
concentration. And short intervals produce thin, one-sided batches,
measured by the unexecuted residual share. Executed volume and trade count
have no theoretical prediction of their own: they check that the two
engines trade the same flow, so that any other difference can be
attributed to the mechanism.

Three adaptations are required for a batch mechanism. Measures that reference
the prevailing midpoint at the moment of the trade have no book to take a
midpoint from during accumulation, so the reference is the clearing price
of the last earlier batch that traded --- one interval $\tau$ old, and
older when recent batches did not trade, never a same-instant quantity as
the CDA's midpoint is. Second, the quoted
spread has no direct analogue during accumulation, so it is defined here
as the gap between the best unfilled buy and the best unfilled sell left
after a batch clears. The third follows from \Cref{sec:fbanotation}: an
FBA trade has no aggressor, so the direction $\varepsilon_k$ in
\eqref{eq:effspread} has to be estimated (\Cref{sec:metrics21}).

Finally, the experimental convention is that of \citet{gode1993allocative}:
hold the induced order flow constant and vary only the mechanism. Their
demonstration that institutions can produce allocative efficiency largely
independently of trader rationality also licenses replay-based simulation,
since much of what a mechanism does is mechanical rather than behavioural;
\citet{lebaron2006agent} surveys the agent-based alternative in which
behaviour is modelled and allowed to adapt, and \citet{wah2013latency}
apply it to latency arbitrage, comparing continuous and periodic call
markets with simulated traders.

\section{From Theory to Design and Measurement}
\label{sec:theorylinks}

\Cref{tab:theorylinks} summarises how each strand of this chapter, and
the boundary argument of the design chapter that builds on them, is used
later: the design block it shapes in \Cref{ch:design}, the metric that
measures it in \Cref{ch:results}, and what this thesis can test. Two
strands, transparency and segregation, inform only the design analysis,
because the replay holds behaviour and the single order pool fixed.

\begin{table}[htbp]
\centering
\caption[From theory to design and measurement]{Each theoretical strand,
its prediction for an FBA, the design block it shapes and the metric that
measures it.}
\label{tab:theorylinks}
\small
\begin{tabularx}{\textwidth}{@{}p{2.6cm} L p{1.9cm} p{3.2cm}@{}}
\toprule
\textbf{Strand} & \textbf{Prediction for the FBA} & \textbf{Design block} & \textbf{Metric (RQ)}\\
\midrule
Latency arbitrage (\ref{sec:discretetime}) & no speed race within $\tau$; lower adverse-selection cost; delay of about $\tau/2$ & 1 Interval & spread decomposition (2.1); time to execution (2.3)\\
Adverse selection (\ref{sec:adverse}) & lower price impact, not zero; price reacts less to flow & 3 Segregation & price impact (2.1); Kyle's lambda (2.2)\\
Transparency (\ref{sec:transparency}) & disclosure trades exploitation against bid shading & 2 Disclosure & not testable in a replay\\
Auction theory (\ref{sec:auctiontheory}) & near-truthful limits; price improvement; thin batches at short $\tau$ & 4 Clearing & trader surplus; residual share (2.3)\\
Market quality (\ref{sec:quality}) & batching pools liquidity and dampens price moves & --- & depth, Amihud (2.1); volatility, dispersion (2.2)\\
Boundary repricing race (\ref{sec:priceimprovement}) & with a fixed, visible close, submissions bunch at the close; speed regains value there & 1, 2, 4 & boundary concentration (2.3), baseline only\\
\bottomrule
\end{tabularx}
\end{table}
